% Faithful conversion of ages-ii-sprint-0; source wording and XML order retained.
\hypertarget{__RefHeading___Toc4785_3118176633}{}
\refstepcounter{subsection}
\subsection*{3.3.1 Sprint 0}
\label{unit:ages-ii-sprint-0}
\addcontentsline{toc}{subsection}{3.3.1 Sprint 0}

% text-p-0844
\indent \originaltab{}Na Sprint 0, iniciei os trabalhos com um foco voltado ao estudo de conceitos básicos de \textit{web scraping}, uma tecnologia que poderia ser útil para a obtenção de dados provenientes das redes sociais. No entanto, à medida que aprofundei os estudos, percebi que essa prática enfrenta diversas barreiras. Muitas plataformas impõem restrições significativas ao \textit{web scraping}, tornando essa estratégia não apenas complexa e de difícil escalabilidade, mas também possivelmente ilegal. Diante disso, optamos por deixar a abordagem em segundo plano.\par

% text-p-0845
\indent \originaltab{}Como nosso projeto não requer uma diversidade de interfaces de usuário, concentrei meus esforços para a pesquisa e exploração de APIs de redes sociais. Isso se mostrou essencial, pois a obtenção de dados de usuários em plataformas como Facebook, Instagram e LinkedIn era uma grande incógnita para o time. Estávamos incertos quanto à viabilidade de acessar tais informações, considerando as restrições impostas pela da Lei Geral de Proteção de Dados (LGPD) e às rigorosas políticas de privacidade de dados das gigantes Meta e Microsoft. Esse questionamento permaneceu em aberto durante a Sprint 0, com respostas mais concretas sendo obtidas apenas na Sprint 1.\par

% text-p-0846
\indent \originaltab{}Durante o desenvolvimento da Sprint, tive de pausar minhas pesquisas para iniciar as discussões sobre o Banco de Dados com meus demais colegas AGES II. Fomos informados pelo cliente, durante a apresentação do projeto, que eles estavam utilizando o banco de dados da Receita Federal e que nos enviariam uma amostra para que pudéssemos trabalhar em cima dela. Prontamente, reuni informações de Bancos de Dados que poderíamos utilizar no projeto e apresentei\originalspace{}\originalspace{}aos AGES II. A escolha do banco é crítica, visto que é o centro de nosso projeto.\originalspace{}\par

% text-p-0847
\indent \originaltab{}Diversos fatores motivaram a escolha do PostgreSQL como nosso Banco de Dados. Inicialmente, a compatibilidade com o sample fornecido pelo cliente facilitou a integração inicial. No entanto, essa não foi a única razão para a nossa escolha. O PostgreSQL se destacou por ser uma plataforma robusta e flexível, especialmente adequada para a manipulação de dados complexos, uma característica essencial para o nosso projeto. Sabíamos que precisaríamos realizar integrações complexas entre os dados obtidos via APIs das redes sociais e o Banco de Dados da Receita Federal, e o PostgreSQL oferecia os recursos necessários para automatizar e manter essas operações de forma eficiente.\par

% text-p-0848
\indent \originaltab{}Como o Banco de Dados precisava ser definido o quanto antes, devido a sua importância para o andamento do projeto, decidi iniciar a modelagem conceitual ainda durante a Sprint 0. Sabíamos que todo o desenvolvimento posterior dependeria dessa modelagem, pois ela determinaria como os dados seriam estruturados e relacionados. No entanto, o cliente só enviou a amostra do banco de dados da Receita Federal no final da Sprint 0, o que limitou a profundidade do trabalho inicial. A modelagem que fiz foi, portanto, preliminar, servindo apenas para esboçar as possíveis relações entre as tabelas. Foi somente na Sprint 1, após a recepção do \textit{database} \textit{sample} e uma análise mais detalhada, que conseguimos realizar a modelagem conceitual e lógica de forma completa e alinhada com as necessidades reais do projeto.\par

\begin{reflection}
% text-p-0851
\indent \originaltab{}A primeira Sprint foi inteiramente dedicada aos estudos, principalmente de APIs e Banco de Dados. Nosso projeto é muito complexo e fora do comum. Não é apenas um \textit{CRUD} (\textit{Create, Read, Update, Delete}) básico. Devido a isso, investi muito tempo estudando a viabilidade de certas ferramentas e estratégias para dar vida ao projeto. Confesso que não colocar a “mão na massa” na Sprint 0 foi um pouco desanimador, mas compreendo que nosso projeto necessita de muita pesquisa para criarmos um software funcional e alcançar as expectativas do cliente.\originalspace{}\par

% text-p-0852
\indent \originaltab{}Durante toda a Sprint 0, fiquei verdadeiramente preocupado se seria possível construirmos uma solução dessa complexidade em apenas 4 meses. Quanto mais eu explorava as APIs de redes sociais, mais incerto ficava. Finalizei a Sprint com muitos questionamentos, mas sabendo que ainda haveria muita pesquisa a ser feita.\par

\end{reflection}
